On the domain depicted in figure~\ref{40000034:domain}, the equation
\begin{equation}
Lu=\lambda u
\qquad\text{with}\qquad
L
=
\frac{\partial^2}{\partial x^2}
-
\frac{\partial}{\partial y}
+
\frac{\partial^2}{\partial y^2}
\label{40000034:pde}
\end{equation}
with homogeneous boundary conditions is given.
Convert the equation~\eqref{40000034:pde} into two ordinary differential
equations and suitable boundary conditions and solve one of them.
\begin{figure}[ht]
\centering
\begin{tikzpicture}[>=latex,thick,scale=4]
\fill[color=red!20] (0,0) -- (1,0) -- +(45:1) -- (45:1) -- cycle;
\fill[color=gray!20,opacity=0.5] (0,0) -- (0.3,0) arc (0:45:0.3) -- cycle;
\node at (22.5:0.2) {$\frac{\pi}4$};
\draw[color=red,line width=1.4pt] (0,0) -- (1,0) -- +(45:1) -- (45:1) -- cycle;
\draw[->] (-0.1,0) -- (1.4,0) coordinate[label={$x$}];
\draw[->] (-135:0.1) -- (45:1.2) coordinate[label={above left:$y$}];
\draw (1,-0.02) -- (1,0.02);
\node at (1,0) [below] {$1$};
\draw ($(45:1)+(-45:0.02)$) -- +(135:0.02);
\node at (45:1) [above left] {$1$};
\node[color=red] at ($0.5*(1,0)+0.5*(45:1)$) {$\Omega$};
\end{tikzpicture}
\caption{Domain for problem~\ref{40000034}
\label{40000034:domain}}
\end{figure}

\begin{loesung}
We write the solution as a product $u(x,y) = X(x)Y(y)$ and substitute
that into the differential equation~\eqref{40000034:pde} giving
\[
X''(x)Y(y) - X(x)Y'(y) + X(x)Y''(y) = \lambda X(x)Y(y)
\qquad\Rightarrow\qquad
\frac{X''(x)}{X(x)}
-
\frac{Y'(y)}{Y(y)}
+
\frac{Y''(x)}{Y(y)}
=
\lambda.
\]
The variables $x$ and $y$ can now be separated to
\[
\frac{X''(x)}{X(x)}
=
\lambda -\frac{Y''(y)}{Y(y} + \frac{Y'(y)}{Y(y)}.
\]
Since the left hand side only depends on $x$ and the right hand side
only depends on $y$, both must be constant.
We call this constant $-\mu^2$ and get the equations
\begin{align*}
X''(x)
&=
-\mu^2 X(x)
&&\text{and}&
Y''(y) - Y'(y) &= (\mu^2-\lambda)Y(y).
\end{align*}
The equation for $X(x)$ has trigonometric functions as solutions,
but only the sine function
\[
X(x) = \sin \mu x
\]
also satisfies the homogeneous boundary condition at $x=0$.
The homogenous boundary condition at $x=1$ implies that $\mu = \pi k$
with $k\in\mathbb{N}$.
\end{loesung}

\begin{diskussion}
The solutions for $Y(y)$ are of the form $Y(y) = e^{ay}\sin by$.
Substituting this into the equation gives
\begin{align*}
Y'(y)  &= ae^{ay}\sin by + b e^{ay} \cos by\\
Y''(y) &= a^2e^{ay}\sin by + abe^{ay}\cos by + bae^{ay}\cos by - b^2e^{ay}\sin by
\\
&= (a^2 -b^2) e^{ay}\sin by + 2abe^{ay} \cos by.
\end{align*}
Substituting this in the 
\begin{align*}
(a^2-b^2-a)e^{ay}\sin by
+(2ab-b)e^{ay}\cos by
&=
(\mu^2+\lambda)e^{ay}\sin by
\\
(a^2-b^2-a-\mu^2-\lambda)e^{ay}\sin by
+
(2a-1)be^{ay} \cos by
&=
0.
\end{align*}
The cosine term can only disappear at $y=0$ if the coefficient is zero,
which implies that $a=-\frac12$. 
Then the first coefficient only disappears if
\[
0
=
a^2-b^2-a-\mu^2-\lambda
=
\frac14-b^2+\frac12-\pi^2k^2-\lambda
=
\frac34 - b^2 -\pi^2k^2-\lambda.
\]
On the other hand, the homogeneous boundary condition at $y=1$ implies that
$b = \pi l$ allowing us to conclude that the possible $\lambda$-values
are
\[
\lambda_{k,l}
=
\frac34
-\pi^2 l^2
-\pi^2 k^2.
\]
Thus the corresponding solutions are
\[
u_{k,l}(x,y)
=
\sin(k\pi x)
\sin(l\pi y)
e^{-y/2}
\]
with $k,l\in\mathbb{N}$.
\end{diskussion}

\begin{bewertung}
Product ansatz ({\bf A}) 1 point,
separation of variables ({\bf S}) 1 point,
two separated differential equations ({\bf D}) 2 points,
solutions of the equations ({\bf S}) 1 point,
restriction of the value of $\mu$ from boundary conditions ({\bf B}) 1 point.
\end{bewertung}

